\documentclass[a4paper,11pt]{article}
\usepackage[utf8]{inputenc}
\usepackage[T1]{fontenc}
\usepackage[french]{babel}
\usepackage[left=2cm,right=2cm,top=2cm,bottom=2cm]{geometry}
\usepackage{xcolor}
\usepackage{hyperref}
\usepackage{fontawesome5}
\usepackage{parskip}

% --- Couleurs Radiall ---
\definecolor{primary}{RGB}{0, 48, 135}
\hypersetup{colorlinks=true, urlcolor=primary}

\begin{document}

% --- EXPÉDITEUR ---
\noindent
\begin{minipage}[t]{0.5\textwidth}
    \textbf{Loïc Aron MBASSI EWOLO}\\
    Élève-Ingénieur Bac+5 -- Double diplôme ENSEA\\[3pt]
    \faMapMarker\ Cergy, France\\
    \faPhone\ (+33) 7 59 52 33 98\\
    \faEnvelope\ \href{mailto:loicaron.mbassiewolo@ensea.fr}{loicaron.mbassiewolo@ensea.fr}
\end{minipage}
\hfill
% --- DATE ---
\begin{minipage}[t]{0.4\textwidth}
    \raggedleft
    Cergy, le 20 février 2026
\end{minipage}

\vspace{1.2cm}

% --- DESTINATAIRE ---
\hfill
\begin{minipage}[t]{0.5\textwidth}
    \raggedleft
    \textbf{Radiall}\\
    Service Recrutement\\
    Saint-Quentin-Fallavier, France
\end{minipage}

\vspace{1cm}

\noindent\textbf{Objet :} Candidature au poste d'Ingénieur Développement Produits (Alternance)\\
\textit{Réf : t52mtturws}

\vspace{0.8cm}

Madame, Monsieur,

Concevoir un banc de test automatisé pour valider des modules de transmission sans contact, puis analyser les résultats pour améliorer le produit : c'est le type de défi technique qui motive ma candidature chez \textbf{Radiall}. En double diplôme d'ingénieur à l'\textbf{ENSEA} (systèmes embarqués, électronique) et diplômé de l'\textbf{École Polytechnique de Yaoundé} (mathématiques appliquées), je dispose des fondamentaux en électronique, en programmation C/Python et en automatisation nécessaires pour contribuer à vos développements produits.

Mon expérience la plus directement transposable à vos besoins est le projet \textit{Harmony Gloves}, réalisé au laboratoire IOTA. Au sein de cette équipe, j'ai participé au prototypage d'un gant instrumenté : \textbf{soudure} de composants sur circuit imprimé, intégration de capteurs IMU et flex sensors, et développement du firmware en C pour l'acquisition de données en temps réel. Ce travail m'a confronté aux contraintes concrètes de la conception hardware, de l'interfaçage capteurs et de la validation par mesures sur banc.

En parallèle, je développe une solide compétence en \textbf{automatisation et scripting Python}. Pour un projet personnel, j'ai conçu un outil de détection vidéo automatisé (YOLOv8/OpenCV) compilé en exécutable distributable, ce qui m'a formé à la création de scripts robustes et à l'optimisation mémoire. Dans le cadre du challenge robotique \textit{CoHoMa} (Défense/ENSEA), je travaille en \textbf{C/C++} sur cible Nvidia Jetson avec intégration Lidar 3D et caméra de profondeur, ce qui renforce ma maîtrise des protocoles de communication et du traitement de données capteurs en environnement embarqué.

Je suis disponible dès \textbf{septembre 2026} pour une alternance de 12 mois, avec la possibilité de débuter par un stage dès \textbf{avril 2026}. Travailler chez Radiall, acteur de référence dans les connecteurs et solutions de transmission, représente pour moi l'opportunité d'appliquer mes compétences en électronique et en automatisation à des produits industriels concrets, tout en progressant au contact de vos équipes R\&D.

Je reste à votre disposition pour un entretien afin de discuter de ma candidature et de vous présenter mes réalisations.

Je vous prie d'agréer, Madame, Monsieur, l'expression de mes salutations distinguées.

\vspace{0.8cm}

\textbf{Loïc Aron MBASSI EWOLO}\\
\textit{Élève-Ingénieur ENSEA / Polytechnique Yaoundé}

\end{document}
